% docs/translation/latex/chapters/05_compat.tex - chapter of the LaTeX edition of the book "translation_report" (SIMC++ documentation).
%   Converted once from docs/translation/markdown/05_compat.md by tools/md2latex.py; from then on THIS FILE is the source of the
%   LaTeX edition and may be edited by hand. Re-running tools/md2latex.py --write-sources for this
%   chapter OVERWRITES the hand edits (see docs/README.md).
%   Compile: cmake --build <build> --target pdf, or tools/build_docs_pdf.py (figures, xelatex passes,
%   makeindex with docs/common/latex/simcpp.ist); by hand, see docs/README.md. XeLaTeX, class
%   docs/common/latex/simcpp.cls (packages: fontspec, tikz, xcolor, hyperref, listings, longtable, booktabs,
%   titlesec, titletoc, fancyhdr, makeidx, preview, ...); fonts TeX Gyre Pagella, TeX Gyre Heros, DejaVu Sans
%   Mono. The figures are built first from docs/{manual,translation}/figures/*.tex (preview.sty) into PDFs.
\chapter{\texorpdfstring{Fortran compatibility layer}{Fortran compatibility layer}}\label{translation_report-05u5fcompat}
\simcabstract{Policy (D-031): \textbf{clear Fortran bugs are corrected by default}; the run option \code{fortran\_\allowbreak{}compat} restores the exact
Fortran behaviour so that the bitwise tests (L1 tables, captures, RANLUX full runs, the compat regression baselines) can
still prove the translation. The complete list of differences is FORTRAN\_\allowbreak{}DIFFERENCES.md\footnote{From SIMC to SIMC++, ``Differences between SIMC++ and the Fortran SIMC'' (\code{docs/\allowbreak{}translation/\allowbreak{}markdown/\allowbreak{}FORTRAN\_\allowbreak{}DIFFERENCES.\allowbreak{}md})}.}
\section{\texorpdfstring{Categories and markers}{Categories and markers}}\label{translation_report-05u5fcompat:51-categories-and-markers}
\begin{simctable}{footnotesize}{5}{>{\raggedright\arraybackslash}p{\dimexpr 0.094\linewidth-2\tabcolsep\relax}>{\raggedright\arraybackslash}p{\dimexpr 0.195\linewidth-2\tabcolsep\relax}>{\raggedright\arraybackslash}p{\dimexpr 0.116\linewidth-2\tabcolsep\relax}>{\raggedright\arraybackslash}p{\dimexpr 0.178\linewidth-2\tabcolsep\relax}>{\raggedright\arraybackslash}p{\dimexpr 0.318\linewidth-2\tabcolsep\relax}}
%
\toprule\rowcolor{simchead}\simcth{Category} & \simcth{meaning} & \simcth{default mode} & \simcth{\code{-{}-{}fortran-{}compat}} & \simcth{marker in the code} \\\midrule\endfirsthead
\multicolumn{5}{@{}l}{\simccontinued}\\\toprule\rowcolor{simchead}\simcth{Category} & \simcth{meaning} & \simcth{default mode} & \simcth{\code{-{}-{}fortran-{}compat}} & \simcth{marker in the code} \\\midrule\endhead
\midrule\multicolumn{5}{r@{}}{\simccontinues}\\\endfoot
\bottomrule\endlastfoot
A & clear Fortran bug & corrected & Fortran & \code{/\allowbreak{}/\ FORTRAN-{}BUG(\allowbreak{}D-{}xxx):\ original\ Fortran\ <file:\allowbreak{}line>:\ <code>\allowbreak{};\ corrected\ because\ <reason>} \\
A-NP & clear bug in a path that is not ported & note only & — & (none left at P15.D1) \\
B & ambiguous, maybe intended (the "expert list") & Fortran & Fortran & \code{/\allowbreak{}/\ FORTRAN-{}DIFF-{}CANDIDATE(\allowbreak{}D-{}xxx):\ .\allowbreak{}.\allowbreak{}.\ Question:\ .\allowbreak{}.\allowbreak{}.} \\
C & open question with the user (acceptance mode, D-009/\allowbreak{}D-028) & Fortran & Fortran & \code{/\allowbreak{}/\ FORTRAN-{}OPEN(\allowbreak{}D-{}031\ flag\ 1)} \\
D & faithful odd behaviour, not a bug (precision, formats, dead code, guards) & Fortran & Fortran & \code{/\allowbreak{}/\ FORTRAN-{}QUIRK(\allowbreak{}.\allowbreak{}.\allowbreak{}.\allowbreak{})} \\
DEV & deliberate design deviation of SIMC++ (JSON input, exit codes, no \code{.\allowbreak{}bin}, arm code 6 refused, ...) & SIMC++ & SIMC++ & FORTRAN\_\allowbreak{}DIFFERENCES.md §6 \\
\end{simctable}
Who decides: findings come from the Fortran reports (\code{reference/\allowbreak{}P*.\allowbreak{}F*\_\allowbreak{}report.\allowbreak{}md}) or from the C++ reading; the director
classifies them in a decision (\code{D-{}0xx}: e.g. D-057 for rho, D-061 SOS/HRS, D-067 target field).

\begin{simcnote}
an agent that suspects a bug keeps the Fortran behaviour and reports it; it does not correct on its own.
\end{simcnote}
\section{\texorpdfstring{How the switch reaches the code}{How the switch reaches the code}}\label{translation_report-05u5fcompat:52-how-the-switch-reaches-the-code}
There is no global. The flag is set once and passed down explicitly:

\begin{simcplain}[basicstyle=\ttfamily\tiny]
simc++ --fortran-compat  |  JSON "run": {"fortran_compat": true}
  → driver::RunOptions::fortran_compat                         (CommandLine.cpp)
  → SetupOptions::fortran_compat → RunSetup::fortran_compat    (BuildRunSetup; set-up corrections use it directly)
  → every stage reads StageContext::setup.fortran_compat in init()
  → kernels get it explicitly from their stage:
        function arguments            SpectrometerAngles(..., fortran_compat), rotateHaxis/rotateVaxis(..., fortran_compat)
        argument structs              ExtremeTripArgs::fortran_compat, proton::DeForestState::fortran_compat,
                                      pion::PeepiConfig::fortran_compat, the Saghai/peeK/peepiX/peerho/peedelta configs
        kernel setters                optics::CosyEngine::setFortranCompat (via McHms/McShms/McSos/McHrs::setFortranCompat),
                                      calo::McCalo::setFortranCompat
        driver                        PassCuts(setup, recon), MonteCarlo, PizeroCalorimeter, RunReport, the sinks
                                      (NtupleLayout, FortranHistSink, GenHistSink: report-level corrections)
\end{simcplain}
The run log prints the mode (\code{\ SIMC++:\ fortran\_\allowbreak{}compat\ =\ true/\allowbreak{}false\ .\allowbreak{}.\allowbreak{}.}), and the summary JSON records it
(\code{run.\allowbreak{}fortran\_\allowbreak{}compat}, D-054), which \code{simc\_\allowbreak{}compare} uses to choose the pass\_cuts window (chapter 2).
\code{-{}-{}fortran-{}compat}\simcidx{Command-line options!fortran-compat@\code{-{}-{}fortran-{}compat}} also refuses \code{-{}-{}threads\ N\ >\ 1}\simcidx{Command-line options!threads@\code{-{}-{}threads}} (the exact reproduction needs one stream in the Fortran order).

Where the switches live (files that read the flag): src/setup (5 files), src/driver (10), src/output (3),
src/\allowbreak{}material,\allowbreak{} src/optics (2), src/\allowbreak{}radiation (2), src/\allowbreak{}reactions (+ pion, rho, sidis), src/\allowbreak{}physics/\allowbreak{}\{proton,\allowbreak{} nuclear,
pion, kaon, delta, sidis, rho\}, src/\allowbreak{}spectrometers/\allowbreak{}\{hms,\allowbreak{} shms, sos, hrs, calo\}. \code{grep\ -{}rn\ fortran\_\allowbreak{}compat\ src\ include}
lists them.

\begin{figure}[htbp]
\simcfigure{compat_switch}
\caption[The two behaviour modes]{The two behaviour modes: what each Category of finding does by default and with -{}-fortran-compat,\allowbreak{} and how the switch reaches the code}\label{fig:compatu5fswitch}
\end{figure}
\section{\texorpdfstring{Adding a correction (Category A)}{Adding a correction (Category A)}}\label{translation_report-05u5fcompat:53-adding-a-correction-category-a}
\begin{enumerate}
\item \textbf{Decision.} The finding must have a D-number and Category A (DECISIONS.md). Note the Fortran file:line, what the
Fortran does, what the correction does, and whether the correction changes the number of random draws (D-058).
\item \textbf{Code.} Keep the Fortran path and add the corrected one behind the flag the kernel already receives (or pass the
flag through its config struct / setter; never read a global):

\simcsource{src/\allowbreak{}driver/\allowbreak{}Simc\allowbreak{}Run.\allowbreak{}cpp}
\begin{simccode}[basicstyle=\ttfamily\tiny,style=simccpp]
// FORTRAN-BUG(D-006): original Fortran simc.f:237: `recon%p%delta .ge. (SPedge%e%delta%max-slop%MC%p%delta%used)`;
// corrected because the hadron-arm delta must be tested against the hadron-arm limit SPedge%p%delta%max, like the
// other five hadron-arm tests (copy-paste of l.231). fortran_compat keeps the electron-arm limit.
bool PassCuts(const RunSetup& S, const Event& r) {
    const auto& E = S.SPedge;
    const auto& M = S.slop.MC;
    const double p_delta_max = S.fortran_compat ? E.e.delta.max : E.p.delta.max;
\end{simccode}
(\code{driver::\allowbreak{}PassCuts}, the first correction of the project.) The compat path must stay bitwise the
Fortran, including its random draws.
\item \textbf{Both-mode test.} Add an item to \code{tests/\allowbreak{}fortran\_\allowbreak{}diff/\allowbreak{}test\_\allowbreak{}fortran\_\allowbreak{}differences.\allowbreak{}cpp} (a function registered in
\code{items(\allowbreak{})} under the id, e.g. \code{"D061\_\allowbreak{}F1e\_\allowbreak{}hrs\_\allowbreak{}vdc\_\allowbreak{}rotation"}) that shows \textbf{both} behaviours: \code{fortran\_\allowbreak{}compat\ =\ true}
gives the Fortran result (ideally against a reference table or run), the default gives the corrected one. Add the id
to \code{SIMC\_\allowbreak{}FORTRAN\_\allowbreak{}DIFF\_\allowbreak{}ITEMS} in tests/\allowbreak{}fortran\_\allowbreak{}diff/\allowbreak{}CMakeLists.\allowbreak{}txt (CTest \code{fortran\_\allowbreak{}diff\_\allowbreak{}<id>}, labels
\code{fortran\_\allowbreak{}diff;unit}). Corrections visible only in a full run can instead be tested by a RANLUX full run in both modes
(tests/l1, \code{ranlux\_\allowbreak{}fullrun.\allowbreak{}sh}).
\item \textbf{The L1 / replay tests keep running in compat} (they pass \code{fortran\_\allowbreak{}compat\ =\ true} or \code{-{}-{}fortran-{}compat}\simcidx{Command-line options!fortran-compat@\code{-{}-{}fortran-{}compat}}); check
that they still pass bitwise.
\item \textbf{FORTRAN\_\allowbreak{}DIFFERENCES.md row} in §2 (Category A): \code{|\ id\ |\ Fortran\ file:\allowbreak{}line\ |\ original\ behaviour\ |\ C++\ default\ |\ effect\ (\allowbreak{}outputs,\ size,\ configurations)\ |\ test\ |}, plus the measured effect in the per-reaction L2 section when the
correction changes distributions (an L2 run in both modes, \code{simc\_\allowbreak{}compare} corrected vs Fortran).
\item \textbf{Baselines.} A correction changes the corrected-mode regression baselines (\code{regression\_\allowbreak{}corrected\_\allowbreak{}*} and the
\code{*\_\allowbreak{}corrected\_\allowbreak{}*} families): regenerate them (chapter 2 §2.6) and state in the commit which families changed and why.
The compat baselines must not change.
\item \textbf{User manual}: one line in the user-facing summary of corrections.
\end{enumerate}
\section{\texorpdfstring{Ambiguous cases, quirks and the acceptance mode}{Ambiguous cases, quirks and the acceptance mode}}\label{translation_report-05u5fcompat:54-ambiguous-cases-quirks-and-the-acceptance-mode}
\begin{itemize}
\item Category B: keep the Fortran in both modes, write \code{FORTRAN-{}DIFF-{}CANDIDATE(\allowbreak{}D-{}xxx)} with the question, add a §4 row in
FORTRAN\_\allowbreak{}DIFFERENCES.md. These items form the expert list (P15.X) for the user and a physicist.
\item Category D: keep, write \code{FORTRAN-{}QUIRK(\allowbreak{}.\allowbreak{}.\allowbreak{}.\allowbreak{})} with the reason (e.g. \code{FORTRAN-{}QUIRK(\allowbreak{}D-{}012\ F2)}), add a §5 row when it
can surprise a user.
\item Category C: the acceptance mode (\code{doing\_\allowbreak{}phsp}\simcidx{Parameters!doing_phsp@\code{doing\_phsp}}, D-009/\allowbreak{}D-028) stays Fortran with \code{FORTRAN-{}OPEN(\allowbreak{}D-{}031\ flag\ 1)} until the
user answers.
\end{itemize}
